\section{UDML Adaptation for P2VG}
\label{sec:udml_p2vg}

This section describes the UDML-style fusion module used to replace the original gated sagittal--axial fusion in P2VG.  The formulation follows the released UDML implementation as closely as possible, while adapting the classification losses of UDML to language-model captioning losses.

\subsection{Notation}

For a training sample, let
\[
    x_s,\; x_a
\]
denote the sagittal and axial MRI inputs, respectively. The two 3D vision encoders produce patch-token features
\[
    H_s = E_s(x_s) \in \mathbb{R}^{B \times N \times C_v},
    \qquad
    H_a = E_a(x_a) \in \mathbb{R}^{B \times N \times C_v},
\]
where \(C_v=768\) for the ViT3D backbone. In MedGemma adapter mode, each view is first spatially pooled and mapped into the native MedGemma visual-projection input space:
\[
    Z_s = A(P(H_s)) \in \mathbb{R}^{B \times N' \times C_m},
    \qquad
    Z_a = A(P(H_a)) \in \mathbb{R}^{B \times N' \times C_m},
\]
where \(P(\cdot)\) is spatial pooling, \(A(\cdot)\) is the learnable adapter, \(N'=256\), and \(C_m=1152\). The UDML fusion module operates on \(Z_s\) and \(Z_a\), not on the original \(768\)-dimensional ViT features.

\subsection{Controlled Noise and Variance Targets}

During training, controlled Gaussian perturbation is optionally applied to each modality:
\[
    \tilde{x}_m
    =
    \operatorname{clip}\left(x_m + \epsilon_m,\;0,\;1\right),
    \qquad
    \epsilon_m \sim \mathcal{N}(0, \eta_m^2 I),
    \qquad
    m \in \{s,a\}.
\]
In the implementation, the sampled noise label is
\[
    \sigma_m \in \{1\} \cup \{\sigma_{\min}, \ldots, \sigma_{\max}\},
\]
and the actual perturbation standard deviation is
\[
    \eta_m = \lambda_{\text{noise}} \sigma_m.
\]
If UDML noise is disabled, or if a sample is selected as clean, then
\[
    \sigma_m = 1.
\]

\subsection{Noise-Aware Variance Estimation}

The token features are pooled to view-level vectors:
\[
    \bar{z}_s = \frac{1}{N'} \sum_{i=1}^{N'} Z_{s,i},
    \qquad
    \bar{z}_a = \frac{1}{N'} \sum_{i=1}^{N'} Z_{a,i}.
\]
The modality uncertainty estimators predict positive scale values:
\[
    \hat{\sigma}_s
    =
    \operatorname{softplus}\left(g_s(\operatorname{stopgrad}(\bar{z}_s))\right)
    + \varepsilon,
    \qquad
    \hat{\sigma}_a
    =
    \operatorname{softplus}\left(g_a(\operatorname{stopgrad}(\bar{z}_a))\right)
    + \varepsilon.
\]
The stop-gradient operation follows the UDML implementation: the variance-estimation loss trains the estimator heads, while preventing this auxiliary objective from directly perturbing the vision encoders.

The variance supervision loss is
\[
    \mathcal{L}_{\text{var}}
    =
    \lambda_{\text{var}}
    \left[
        \left\lVert \hat{\sigma}_s - \sigma_s \right\rVert_2^2
        +
        \left\lVert \hat{\sigma}_a - \sigma_a \right\rVert_2^2
    \right].
\]

\subsection{Uncertainty-Based Target Weights}

Following the released UDML weighting code, the inverse-uncertainty target weights are computed as
\[
    \tilde{w}_s
    =
    \frac{2\hat{\sigma}_a^2}
         {\hat{\sigma}_s^2+\hat{\sigma}_a^2+\varepsilon},
    \qquad
    \tilde{w}_a
    =
    \frac{2\hat{\sigma}_s^2}
         {\hat{\sigma}_s^2+\hat{\sigma}_a^2+\varepsilon}.
\]
Thus, if the sagittal uncertainty \(\hat{\sigma}_s\) is high, the sagittal target weight \(\tilde{w}_s\) is reduced, and vice versa.

\subsection{Shared-Weight Dependency Calculator}

To mimic the share-weight auxiliary fusion in UDML, a single shared head \(q(\cdot)\) is applied to full and unimodal feature combinations:
\[
    r_{\text{full}}
    =
    q([\bar{z}_s;\bar{z}_a]),
\]
\[
    r_s
    =
    q([\bar{z}_s;\mathbf{0}]),
    \qquad
    r_a
    =
    q([\mathbf{0};\bar{z}_a]).
\]
The same parameters of \(q\) are reused for all three branches, analogous to the shared \(fc\_out\) in UDML's auxiliary fusion module.

The dependency scores are computed from the unimodal auxiliary responses:
\[
    \alpha_s
    =
    \sum_j
    \left|
        \frac{1}{B}\sum_{b=1}^{B} r_{s,bj}
    \right|
    + \varepsilon,
    \qquad
    \alpha_a
    =
    \sum_j
    \left|
        \frac{1}{B}\sum_{b=1}^{B} r_{a,bj}
    \right|
    + \varepsilon.
\]
This follows the released UDML code, where modality dependency is estimated from the magnitude of the corresponding unimodal auxiliary output.

\subsection{Unbiased Dynamic Fusion}

The uncertainty weights are corrected by the dependency scores:
\[
    u_s = \frac{\tilde{w}_s}{\alpha_s},
    \qquad
    u_a = \frac{\tilde{w}_a}{\alpha_a}.
\]
They are then normalized to sum to \(2\), matching the two-modality normalization in UDML:
\[
    w_s
    =
    \frac{2u_s}{u_s+u_a+\varepsilon},
    \qquad
    w_a
    =
    \frac{2u_a}{u_s+u_a+\varepsilon}.
\]
The fused MedGemma-space visual tokens are
\[
    Z_f
    =
    \frac{1}{2}
    \left(
        w_s Z_s + w_a Z_a
    \right).
\]
Since \(w_s+w_a \approx 2\), the factor \(1/2\) preserves the feature scale.

The fused tokens are then passed through the pretrained MedGemma visual projection:
\[
    V_f
    =
    W_{\text{MG}}
    \left(
        \operatorname{RMSNorm}(Z_f)
    \right)
    \in \mathbb{R}^{B \times N' \times C_{\text{LLM}}},
\]
where \(C_{\text{LLM}}=2560\) for the MedGemma 4B text backbone.

\subsection{Captioning Losses}

Let \(y=(y_1,\ldots,y_T)\) be the target report tokens. The main language-model loss is
\[
    \mathcal{L}_{\text{full}}
    =
    -\sum_{t=1}^{T}
    \log p_{\theta}
    \left(
        y_t
        \mid
        y_{<t}, V_f
    \right).
\]
To adapt UDML's unimodal classification losses to report generation, sagittal-only and axial-only visual tokens are also projected through the same MedGemma projector and decoded by the same LLM:
\[
    V_s
    =
    W_{\text{MG}}
    \left(
        \operatorname{RMSNorm}(Z_s)
    \right),
    \qquad
    V_a
    =
    W_{\text{MG}}
    \left(
        \operatorname{RMSNorm}(Z_a)
    \right).
\]
The auxiliary unimodal language-model losses are
\[
    \mathcal{L}_{s}
    =
    -\sum_{t=1}^{T}
    \log p_{\theta}
    \left(
        y_t
        \mid
        y_{<t}, V_s
    \right),
\]
\[
    \mathcal{L}_{a}
    =
    -\sum_{t=1}^{T}
    \log p_{\theta}
    \left(
        y_t
        \mid
        y_{<t}, V_a
    \right).
\]
These two losses correspond to UDML's unimodal auxiliary losses, but with causal language-model cross entropy replacing classification cross entropy.

\subsection{Total Objective}

The full training objective is
\[
    \mathcal{L}_{\text{total}}
    =
    \mathcal{L}_{\text{full}}
    +
    \mathcal{L}_{\text{var}}
    +
    \lambda_{\text{uni}}
    \left(
        \mathcal{L}_{s}
        +
        \mathcal{L}_{a}
    \right).
\]
If unimodal LM supervision is disabled, then \(\lambda_{\text{uni}}=0\). If UDML noise supervision is disabled, the variance targets default to the clean value \(\sigma_s=\sigma_a=1\).

\subsection{Relationship to UDML}

The implementation preserves the main mechanisms of UDML:
\[
    \text{controlled noise}
    \;+\;
    \text{variance estimation}
    \;+\;
    \text{share-weight dependency estimation}
    \;+\;
    \text{dependency-corrected dynamic fusion}
    \;+\;
    \text{full and unimodal task losses}.
\]
The task-specific adaptation is that UDML's classification cross entropy is replaced by causal language-model cross entropy for report generation.
